\section{Techno-economic analysis}\label{sec:tea}
\subsection{Economic-model boundary and formulation}
After the carbon ledger establishes the denominator for abatement cost, the economic model constructs the numerator. The principal mature case uses the coherent Nishihara reference economics: 59.7 billion JPY plant cost, 0.52 JPY/MJ heat and 4.9 JPY/kWh electricity at 85\% availability~\cite{nishihara2007potential}. Because no validated project off-design electricity output is claimed, the controlling economic test assigns \textbf{S\$0/MWh electricity export value}. To avoid free reactor capacity, the calculation charges the full published source-product burden: the source 370 MWth heat product at its source heat cost plus the source 88 MWe generation at its source electricity cost. This is a conservative economic allocation, not a claim that the project itself exports 88 MWe.

The source's 70.9 billion JPY / 0.57 JPY/MJ / 5.5 JPY/kWh case is retained only as an adverse \emph{cost} sensitivity in which assigned IHX/secondary-loop cost doubles; it is not interpreted as increased exchanger capacity. CCS capital is anchored to IEAGHG Case 1A and T\&S remains a Singapore screening scenario~\cite{ieaghg2017smr,mticcs2025cost}.

Forward abatement cost remains
\begin{equation}\label{eq:abatement-cost-definition}
C_{\rm abatement}=\frac{C_{\rm candidate}-C_{\rm baseline}}{A_{\rm avoided}}.
\end{equation}
where $C_{\rm candidate}$ and $C_{\rm baseline}$ are annual candidate and baseline costs (S\$/y), $A_{\rm avoided}$ is annual lifecycle abatement (tCO$_2$e/y), and $C_{\rm abatement}$ is forward abatement cost (S\$/tCO$_2$e). Section~\ref{sec:ledgers} provides the numerical substitution.



\subsection{Historical one-train differential screen (superseded)}
The following cost bridge is retained for numerical provenance, not as the controlling CN4252 economic answer. It predates the stronger E1 full-burden accounting and the E2B two-train maturation model. Its S\$3.725/t result must not be interpreted as a deployable FOAK price.\n\n\subsection{Techno-economic calculation}
The natural-gas saving is directly reproducible from the same cited 52.5 and 34.0 MMSCFD source rates. Assumption A-08 sets the gas price $P_{\rm NG}$ to S\$15/GJ as a project screening value, not a contracted Singapore tariff. Annual gas expenditure is
\begin{equation}\label{eq:ng-cost-general}
C_{\rm NG}=Q_{\rm NG,d}(365)A P_{\rm NG},
\end{equation}
where $C_{\rm NG}$ is annual gas cost (S\$/y). For the cited baseline flow~\cite{inltev953},
\begin{equation}\label{eq:ng-cost-base}
C_{\rm NG,base}
=\left[(52.5)(10^6)(1044)(1.05505585262\times10^{-6})\right]
(365)(0.85)(15)
\approx\mathrm{S\$}269.115~\mathrm{million/y},
\end{equation}
where the bracketed term is GJ/d and Assumption A-08 is $P_{\rm NG}=\mathrm{S\$}15$/GJ. For the cited candidate flow~\cite{inltev961},
\begin{equation}\label{eq:ng-cost-candidate}
C_{\rm NG,cand}
=\left[(34.0)(10^6)(1044)(1.05505585262\times10^{-6})\right]
(365)(0.85)(15)
\approx\mathrm{S\$}174.284~\mathrm{million/y}.
\end{equation}
The factors have dimensions MMSCFD$\times$scf/MMSCF$\times$Btu/scf$\times$GJ/Btu$\times$d/y$\times$S\$/GJ, leaving S\$/y.
hence
\begin{equation}\label{eq:ng-saving}
S_{\rm NG}=269.115-174.284
=\mathbf{S\$94.831~million/y}.
\end{equation}

The selected source economic burden is deliberately charged in full. Nishihara reports source heat cost $p_{h,0}=0.52$ JPY/MJ and source electricity cost $p_{e,0}=4.9$ JPY/kWh~\cite{nishihara2007potential}. Let $D_{JP}=112.27/99.59=1.12732$ be the recorded 2007-to-2025 Japan GDP-deflator factor (-) and $X_{JPY}=0.008117$ SGD/JPY the dated project FX conversion. The normalized heat price is
\begin{equation}\label{eq:heat-price-normalization}
\begin{aligned}
p_h&=p_{h,0}D_{JP}X_{JPY}(1000~\mathrm{MJ/GJ})\\
&=(0.52~\mathrm{JPY/MJ})(1.12732)(0.008117~\mathrm{SGD/JPY})\\
&\qquad\times(1000~\mathrm{MJ/GJ})
\approx\mathrm{S\$}4.758/\mathrm{GJ}.
\end{aligned}
\end{equation}
Similarly,
\begin{equation}\label{eq:electric-price-normalization}
\begin{aligned}
p_e&=p_{e,0}D_{JP}X_{JPY}(1000~\mathrm{kWh/MWh})\\
&=(4.9~\mathrm{JPY/kWh})(1.12732)(0.008117~\mathrm{SGD/JPY})\\
&\qquad\times(1000~\mathrm{kWh/MWh})
\approx\mathrm{S\$}44.837/\mathrm{MWh}.
\end{aligned}
\end{equation}
Here $p_h$ and $p_e$ are normalized source-product prices. The deflator and FX bases are documented project normalization inputs; the GDP deflator is a broad screening index, not a nuclear-construction index. Therefore
\begin{equation}\label{eq:source-heat-cost}
C_{\rm source,heat}
=(370~\mathrm{MW_{th}})(7446~\mathrm{h/y})(3.6~\mathrm{GJ/MWh})(4.758~\mathrm{S\$/GJ})
\approx\mathrm{S\$}47.193~\mathrm{million/y}.
\end{equation}
The 370 MWth source process-heat/economic branch is from Nishihara~\cite{nishihara2007potential}; 7446 h/y is $8760\times0.85$ using Assumption A-02; 3.6 GJ/MWh is an exact energy conversion; and 4.758 S\$/GJ is Eq.~\ref{eq:heat-price-normalization}. The dimensions reduce as MW$\times$h/y$\times$GJ/MWh$\times$S\$/GJ = S\$/y. This is the heat-product share of the selected published source economic burden.
\begin{equation}\label{eq:source-electric-cost}
C_{\rm source,electric}
=(88~\mathrm{MWe})(7446~\mathrm{h/y})(44.837~\mathrm{S\$/MWh})
\approx\mathrm{S\$}29.380~\mathrm{million/y}.
\end{equation}
The 88 MWe complementary source product is the rounded Nishihara high-hydrogen economic configuration~\cite{nishihara2007potential}; 44.837 S\$/MWh is Eq.~\ref{eq:electric-price-normalization}. MW$\times$h/y$\times$S\$/MWh reduces to S\$/y.
and
\begin{equation}\label{eq:source-total-cost}
C_{\rm source}\approx\mathbf{S\$76.572~million/y}.
\end{equation}
The 88 MWe term is used only to recover the full published source economic burden; it is not a project export claim.

IEAGHG Case 1A provides the source capture-capital anchor $C_{\rm CCS,0}=\mathrm{EUR}\,41.02$ million (2014 basis)~\cite{ieaghg2017smr}. The canonical reference capture scale is $M_0=387{,}805.2$ t/y and the project capture scale is $M_{\rm captured}=542{,}361.98$ t/y from Eq.~\ref{eq:captured-annual}. The project assumes linear captured-tonne scaling, exponent $\alpha=1$, together with the recorded 2014-to-2025 deflator proxy $D_{EU}=123.84/90.46=1.36834$ and dated FX $X_{EUR}=1.452$ SGD/EUR. Thus
\begin{equation}\label{eq:ccs-capex-general}
C_{\rm CCS}=C_{\rm CCS,0}D_{EU}X_{EUR}
\left(\frac{M_{\rm captured}}{M_0}\right)^{\alpha},
\end{equation}
where $C_{\rm CCS}$ is project screening capture CAPEX (S\$), $C_{\rm CCS,0}$ is source CAPEX (EUR), $D_{EU}$ is the price-year factor (-), $X_{EUR}$ is FX (SGD/EUR), and $\alpha=1$ is the explicit project scaling assumption. Numerically,
\begin{equation}\label{eq:ccs-capex-substitution}
C_{\rm CCS}
=(41.02~\mathrm{million~EUR})(1.36834)(1.452~\mathrm{SGD/EUR})
\left(\frac{542{,}361.98~\mathrm{t/y}}{387{,}805.2~\mathrm{t/y}}\right)
\approx\mathrm{S\$}114.036~\mathrm{million}.
\end{equation}
IEAGHG did not report S\$114.036 million for this project; it is the project-normalized, linearly scaled result. Capital is annualized using
\begin{equation}\label{eq:crf}
\mathrm{CRF}(i,n)=\frac{i(1+i)^n}{(1+i)^n-1},
\end{equation}
where $i$ is annual discount rate (-) and $n$ is economic life (y). Assumption A-13 retains the IEAGHG $i=0.08$ and $n=25$ y basis~\cite{ieaghg2017smr}; therefore
\begin{equation}\label{eq:ccs-annual-cost}
C_{\rm CCS,annual}=C_{\rm CCS}\,\mathrm{CRF}(0.08,25)
=(114.036~\mathrm{million~S\$})(0.09368)
\approx\mathbf{S\$10.683~million/y}.
\end{equation}
$C_{\rm CCS,annual}$ is annualized capture capital (S\$/y), and $\mathrm{CRF}(0.08,25)\approx0.09368$ y$^{-1}$.
Assumption A-14 is an explicit 10\% project integration/site capital allowance rather than an observed Singapore cost. The S\$546.283 million reactor/source capital basis is the normalized 59.7 billion JPY Nishihara source plant cost~\cite{nishihara2007potential}; S\$114.036 million is the project CCS CAPEX from Eq.~\ref{eq:ccs-capex-substitution}. Assumption A-15 applies the retained 3\% annualisation rate and 40-y life to this allowance; these are screening/design-study bases, not a financing offer. It contributes approximately
\begin{equation}\label{eq:integration-cost}
C_{\rm integration}=0.10(546.283+114.036)\,\mathrm{million}\,\mathrm{CRF}(0.03,40)\approx\mathbf{S\$2.857~million/y}.
\end{equation}
while Assumption A-16 sets T\&S to S\$15/tCO$_2$ as a project screening tariff, not an observed or contracted Singapore tariff. Official Singapore evidence is cited only for the separate statement that actual cross-border CCS costs remain unresolved~\cite{mticcs2025cost}. Using $M_{\rm captured}=542{,}362$ t/y from Eq.~\ref{eq:captured-annual},
\begin{equation}\label{eq:ts-cost}
C_{\rm T\&S}=542{,}362(15)
\approx\mathbf{S\$8.135~million/y}.
\end{equation}
Assumption A-17 sets project electricity revenue to exactly \textbf{S\$0/y} because no validated off-design export model exists; this is a conservative model boundary, not a market-price claim.
The economic ledger follows the same logic. The conventional baseline's natural-gas expenditure is replaced by lower candidate gas use plus the full selected source reactor economic burden, CCS annualisation, integration allowance and T\&S. Project electricity revenue is exactly zero.

\begin{table}[htbp]\centering\scriptsize
\caption{Canonical annual-cost ledger generated by Rust for the zero-electricity-credit model boundary.}
\label{tab:cost-ledger}
\pgfplotstabletypeset[col sep=comma,
columns={item,direction,annual_sgd,meaning},
columns/item/.style={string type,column name=Item,column type={p{0.31\linewidth}}},
columns/direction/.style={string type,column name=Direction,column type={p{0.11\linewidth}}},
columns/annual_sgd/.style={fixed,precision=0,column name={S\$/y},column type=r},
columns/meaning/.style={string type,column name=Meaning,column type={p{0.31\linewidth}}},
every head row/.style={before row=\toprule,after row=\midrule},
every last row/.style={after row=\bottomrule}]
{../results/final_design/generated/cost_ledger.csv}
\end{table}

Total money saved is approximately S\$94.831 million/y.  Total new represented costs are approximately S\$98.247 million/y, so
\begin{equation}\label{eq:annual-incremental-cost}
\Delta C_{\rm annual}
\approx98.247-94.831
=\mathbf{S\$3.416~million/y}.
\end{equation}
The final division is therefore
\begin{equation}\label{eq:abatement-cost-final}
C_{\rm abatement}
=\frac{\Delta C_{\rm annual}}{A_{\rm lifecycle}}
\approx\frac{\mathrm{S\$}3.416\ \mathrm{million/y}}{917{,}139\ \mathrm{tCO_2e/y}}
\approx\mathbf{S\$3.725/\mathrm{tCO_2e}}.
\end{equation}

\begin{figure}[htbp]\centering
\input{../results/final_design/generated/cost_bridge_figure.tex}
\caption{Rust-generated annual cost-contribution bridge.  Natural-gas savings are negative cost; reactor, CCS, integration and transport/storage are added costs.  Electricity revenue is zero by model boundary.}
\label{fig:cost-bridge}
\end{figure}




For the preferred mature case, the represented annual ledger is
\[
C_{\rm base,annual}=538.230~\mathrm{S\$\,million/y},\qquad
C_{\rm cand,annual}=616.814~\mathrm{S\$\,million/y},
\]
so that
\[
\Delta C_{\rm annual}
=C_{\rm cand,annual}-C_{\rm base,annual}
\approx78.583~\mathrm{S\$\,million/y}.
\]
The candidate ledger contains the normalized/annualized mature nuclear burden (S\$207.763 million/y), two-train integration/IHX allowance (S\$22.846 million/y), candidate NG, CCS and T\&S terms; the baseline ledger contains the corresponding two-train conventional NG/process cost basis. The project NG assumption remains S\$15/GJ and T\&S remains S\$15/t captured CO$_2$; neither is a contracted tariff. With E2B lifecycle avoidance
\[
A_{\rm lifecycle}=1.834278~\mathrm{MtCO_2e/y},
\]
the abatement-cost definition is
\begin{equation}
C_{\rm abat}
=\frac{C_{\rm cand,annual}-C_{\rm base,annual}}{A_{\rm lifecycle}},
\end{equation}
where $C_{\rm abat}$ is in S\$/tCO$_2$e when annual costs are expressed in S\$ million/y and $A_{\rm lifecycle}$ in MtCO$_2$e/y. Numerically,
\[
C_{\rm abat,mature}
=\frac{616.814-538.230}{1.834278}
\approx\mathbf{42.84~\mathrm{S\$/tCO_2e}}.
\]
Thus the forward mature model is below the S\$100/t threshold, but only for the \textbf{projected/modelled 10-OAK case}. It is not an observed Singapore quotation, contains no electricity or fictional heat-sharing revenue, and remains exposed to unresolved project costs under the E7 gates.

For comparison, the same forward ledger method gives S\$137.74/tCO$_2$e at FOAK and S\$74.14/tCO$_2$e on the medium BOAK basis. The passing future values therefore come from source-supported two-train utilisation plus projected maturation, not from back-solving the learning rate to the assignment threshold.

\subsection{What is not fully priced here?}
The historical one-train differential screening cost is not a turnkey Singapore project price. The current boundary does not fully price project-specific FOAK premium, financing structure, Singapore site/construction premium, licensing and security, detailed EPC/contingency, emergency-planning infrastructure, detailed 176.8 MWth IHX qualification, waste/decommissioning, insurance/liability, schedule risk, or negotiated cross-border CCS infrastructure/contracts. No monetary values are invented for these unresolved items. Their omission is why the economic conclusion is threshold compliance within the declared model, not bankability.

This distinction also strengthens the economics. The controlling S\$0/MWh case does not require electricity revenue. Instead, it charges the full selected source-product economic burden while giving no value to uncertain project power output. The doubled IHX/secondary-loop case remains only an adverse source cost sensitivity, not a capacity argument.

\subsection{Controlling deployment-stage economics}
\label{sec:stage-economics}
The preceding S\$3.725/tCO$_2$e calculation is retained as a \textbf{historical differential screening bridge}; it no longer controls the final economic conclusion. A stronger E1 accounting assigns the full normalized one-module nuclear burden without unsupported co-product sharing and gives approximately \textbf{S\$154/tCO$_2$e}, which fails the S\$100/t threshold.

E2B changes the physical utilisation architecture before applying an independently selected maturation method. INL/RPT-23-72972 supplies the medium non-electric advanced-reactor/HTGR method: a BOAK overnight-cost basis of USD6,000/kWe (2019 USD) and USD25/MWh OPEX is converted to a heat-only basis by removing the 9\% energy-conversion-system share and applying a 40\% HTGR electric-to-thermal conversion; the resulting BOAK heat-only bases are USD2,184/kWt and USD10/MWth-h~\cite{inl2023advancedcost}. FOAK applies the source medium premium of 1.6. Replication follows
\begin{equation}
C_N=C_1(1-LR)^{\log_2 N},
\end{equation}
where $C_N$ is overnight cost at the $N$th replicated unit, $C_1$ is the BOAK heat-only starting cost, $LR=0.10$ is the source medium learning rate and $N$ is cumulative unit number~\cite{inl2023advancedcost}. At $N=10$, this gives approximately USD1,539/kWt (2019 USD); OPEX is held at USD10/MWth-h as in the source example. These are modelled/projected maturation inputs, not observed HTGR construction costs, and they were selected independently of the CN4252 S\$100/t threshold.

The project then applies its documented 2019-to-2025 broad price escalation and SGD/USD conversion before annualising capital and adding O\&M. One 600 MWth source serves two real 176.8 MWth process trains, for 353.6 MWth useful reformer duty. No electricity revenue and no fictional heat-sharing customer are credited. The resulting stage-specific screening economics are:
\begin{table}[htbp]\centering\small
\caption{Controlling deployment-stage economics for the preferred two-train nuclear-assisted architecture. Values are project-model outputs, not vendor quotations.}
\label{tab:deployment-economics}
\begin{tabular}{p{0.31\linewidth}p{0.22\linewidth}p{0.34\linewidth}}
\toprule
Deployment state & Abatement cost & Interpretation \\ \midrule
FOAK two-train & S\$137.74/tCO$_2$e & FAIL; adverse current/FOAK economic evidence \\
Early-commercial/BOAK & S\$74.14/tCO$_2$e & projected/modelled numerical PASS; unresolved-cost margin remains \\
Preferred mature 10-OAK & S\$42.84/tCO$_2$e & projected/modelled future PASS; not observed commercial cost \\
\bottomrule
\end{tabular}
\end{table}
The mature pass arises from \emph{real higher utilisation of the source-supported process-heat branch plus evidence-backed projected maturation}; the same mature cost basis applied to an under-utilised one-train architecture remains above S\$100/tCO$_2$e. Project-specific integration, site, licensing, security, waste, financing and contracted CCS costs remain unresolved and must be incorporated as they become known.
